\clearpage
\chapter{旋律写作}
\label{sec:melody-writing}

\subsection*{学完这一章能做什么}
\begin{itemize}
  \item 从一段旋律里认出动机，分清“动机”与“整句”不是一回事。
  \item 用重复、模进、倒影、逆行与扩缩，把一个动机长成四小节乐句。
  \item 写出有问有答的四小节：问句停在属或非主音，答句回到主音。
  \item 让旋律音与和弦音、张力音各就各位，并在夜半的钢琴卷帘里把它铺出来。
\end{itemize}

\subsection{动机：一首歌的最小记忆单位}

动机（motive）是一段很短、能被记住并反复使用的音乐材料，通常只有两到五个音，或一个鲜明的节奏型。它不需要完整的意思，但必须有可辨认的“形状”：音高轮廓、节奏、重音、音区，至少有一项足够突出。

比如 C–D–E–G 这个上行动机：轮廓是“低—中—高—更高”，节奏可以是“长—短—短—长”。下次听到 D–E–F–A，你会觉得“还是那句话，只是整体挪高了”；听到 C–B–A–F，你会觉得“方向反了，但骨架还在”。这就是动机的识别力。

动机不是“随便几个音”。如果四个音之间没有可重复的轮廓或节奏，它很难在下一句被认出来。写动机时先问自己：我最想让人记住的是哪一个跳进、哪一个节奏，还是哪一个音色？把那一项保住，其他都可以变。

\subsection{乐句与问答：四小节怎么成立}

乐句（phrase）是能独立成句的一段旋律，流行写作里常是四小节或八小节。四小节最常用的结构是“两小节问、两小节答”：

\begin{itemize}
  \item \textbf{问句}（antecedent）：把动机抛出来，句尾停在属功能、非主音或一个未解决的高音，让耳朵等下文。
  \item \textbf{答句}（consequent）：接住问句的动机，句尾回到主音或主功能，给耳朵一个落地。
\end{itemize}

问句不一定要写成“疑问语气”。它只是没有闭合：可能停在 G4（属音）、E4（三音）或一个张力音上；答句也不一定要完全重复问句，它可以换尾、换节奏、换音区，只要把“没说完”变成“说完了”。

常见做法是让问句与答句共用前两小节的动机，只在后两小节分叉。这样既统一又有回答感：耳朵先认出“又是它”，再听到“这次落地了”。如果问句和答句完全一样，通常只是重复，不是问答；如果答句完全换新材料，又会像另一句话。


\begin{center}
\begin{minipage}{0.92\linewidth}
\centering
\begin{tikzpicture}[x=2.55cm,y=1.0cm,font=\scriptsize]
  \draw[->,gray] (-0.2,-0.2) -- (4.35,-0.2) node[right,font=\tiny]{小节};
  \foreach \x/\lab in {0/1,1/2,2/3,3/4} {
    \draw[gray!45] (\x,-0.2) -- (\x,2.9);
    \node[below,font=\tiny] at (\x,-0.2) {\lab};
  }
  \draw[very thick] (0,0.2) -- (1,1.1) -- (2,2.6) -- (3,1.4) -- (4,0.1);
  \foreach \x/\y/\n in {0/0.2/C4,2/2.6/G4,4/0.1/C4} {
    \fill (\x,\y) circle (1.7pt);
    \node[above=1pt,font=\tiny] at (\x,\y) {\n};
  }
  \draw[<->] (0,3.05) -- (2,3.05) node[midway,above,font=\tiny]{问句};
  \draw[<->] (2,3.05) -- (4,3.05) node[midway,above,font=\tiny]{答句};
\end{tikzpicture}

图 7-2：四小节问答的张力与落点。问句在第 2 小节升到 G4，答句在第 4 小节回 C4；曲线只表示张力高低，不是精确音高。
\end{minipage}
\end{center}

\subsection{重复、模进与变奏：让动机长大}

动机要长大，不能只靠原样重复。下面六种手法最常用：

\begin{center}
\small
\begin{tabular}{p{0.13\linewidth}p{0.28\linewidth}p{0.44\linewidth}}
\toprule
手法 & 做法 & 效果 \\
\midrule
重复 & 原样再来一次 & 强调，建立记忆 \\
模进 & 整体挪到另一个音高，轮廓与节奏不变 & 有推进感，像“同一句话换个人说” \\
倒影 & 以某个音为轴，上下方向反过来 & 熟悉又陌生，适合转折 \\
逆行 & 从最后一个音倒着弹回去 & 制造“回顾”或收束感 \\
扩缩 & 把音值整体加倍或减半 & 同一动机变慢/变快，情绪跟着变 \\
加花 & 在骨干音之间加经过音、辅助音 & 更流畅，但别把骨架盖住 \\
\bottomrule
\end{tabular}

表 7-1：六种动机变形。一次只用一两种，堆太多会听不出原动机。
\end{center}

模进最容易被滥用：连续三四次上行模进，会像爬楼梯，听众很快疲劳。通常一到两次模进就够，第三次要换方向、换节奏或干脆收住。倒影与逆行也不要为了“高级”而用；它们和原动机的轮廓差异越大，越需要前后文把关系说清楚。

\begin{center}
\begin{minipage}{0.86\linewidth}
\centering
\begin{tikzpicture}[x=1.55cm,y=0.55cm,xshift=0.55cm,font=\scriptsize]
  % 原型：C D E G（相对高度 0,1,2,4）
  \foreach \x/\y/\n in {0/6/C,1/6.9/D,2/7.8/E,3/9.4/G} {
    \fill (\x,\y) circle (1.6pt);
    \node[above=1pt] at (\x,\y) {\n};
  }
  \draw (0,6) -- (1,6.9) -- (2,7.8) -- (3,9.4);
  
  % 模进：D E F A（整体上行二度）
  \foreach \x/\y/\n in {0/3/D,1/3.9/E,2/4.8/F,3/6.4/A} {
    \fill (\x,\y) circle (1.6pt);
    \node[above=1pt] at (\x,\y) {\n};
  }
  \draw (0,3) -- (1,3.9) -- (2,4.8) -- (3,6.4);
  
  % 倒影：C B A F（以 C 为轴向下）
  \foreach \x/\y/\n in {0/0/C,1/-0.9/B,2/-1.8/A,3/-3.4/F} {
    \fill (\x,\y) circle (1.6pt);
    \node[below=1pt] at (\x,\y) {\n};
  }
  \draw (0,0) -- (1,-0.9) -- (2,-1.8) -- (3,-3.4);
  \end{tikzpicture}

图 7-1：同一动机的三种轮廓。自上而下是原型、模进、倒影：原型上行，模进整体挪高，倒影把方向反过来；图只画轮廓，不对应具体谱表。
\end{minipage}
\end{center}

\subsection{旋律与和弦：什么时候踩和弦音，什么时候用张力音}

旋律不是“把和弦音一个个弹出来”。它要在稳定与不稳定之间走：强拍踩和弦音，弱拍用经过音、辅助音把路走顺，偶尔用张力音制造色彩。

\begin{center}
\begin{minipage}{\linewidth}
\centering
\small
\begin{tabular}{p{0.16\linewidth}p{0.30\linewidth}p{0.40\linewidth}}
\toprule
类型 & C 和弦上的例子 & 用法 \\
\midrule
和弦音 & C、E、G & 强拍、句尾、落点；稳定 \\
经过音 & D（C 与 E 之间） & 弱拍，顺着音阶走到下一个和弦音 \\
辅助音 & F（E 的上方） & 离开 E 再回到 E，制造小波动 \\
倚音 & 强拍上的 F，随后解决到 E & 先用张力抓住耳朵，再落下 \\
张力音 & A（6）、D（9）、B（大七） & 色彩与悬念；要有前后文接住 \\
\bottomrule
\end{tabular}

表 7-2：旋律音与 C 大三和弦的关系。张力音不是“错音”，但需要解决或风格支撑。
\end{minipage}
\end{center}

最稳的顺序是“先写旋律，再配和弦”。如果先定和弦，再让旋律去填每个和弦音，旋律很容易变成琶音练习。实际写作可以先哼出四小节，再问：强拍落在哪个和弦音上？句尾停在主音还是属音？如果旋律的强拍一直撞在张力音上，要么改旋律，要么改和弦，别假装没听见。

\textbf{和弦轨的边界：}夜半的和弦轨是固定五块演示数据，不跟随工程里的和弦；它只能用来听 i–VI–iv–V 的参考，不能当作“当前工程和弦”来给旋律配和声。配和声仍要在卷帘里自己写和弦，或用耳朵判断。

\subsection{节奏动机：同样的音高，换节奏就是另一句话}

节奏动机（rhythmic motive）是不依赖具体音高的记忆点。把 C–D–E–G 的四个音改成“长—短—短—长”，和改成“短—短—短—长”，听众会认为是两句不同的音乐，哪怕音高完全一样。

\begin{center}
\small
\begin{tabular}{p{0.13\linewidth}p{0.30\linewidth}p{0.42\linewidth}}
\toprule
节奏型 & 4/4 里的写法 & 效果 \\
\midrule
均分 & 每拍一个音，或每半拍一个音 & 平稳、推进，适合陈述 \\
附点 & 长音＋短音，如“长—短” & 有弹性，带一点摇摆 \\
切分 & 重音落在弱拍或反拍 & 推力强，流行/放克里常用 \\
三连音 & 一拍里三个音 & 流动、密集，适合过渡或加花 \\
\bottomrule
\end{tabular}

表 7-3：四种基础节奏动机。先用手打拍子确认，再写进卷帘。
\end{center}

写节奏动机时，先用嘴或手打出重音，再决定哪些音落在强拍。切分不是“随便把音挪后半拍”：它要有前后的稳定拍作参照，否则听起来只是拖拍。三连音也不宜连续太多小节，稠密之后需要回到均分，让耳朵休息。

\subsection{在夜半里怎么做}

\begin{enumerate}
  \item 新建 4/4 工程，选中一条 MIDI 乐器轨，打开钢琴卷帘（P/F7）。
  \item 在第 1 小节写四拍动机：C4、D4、E4、G4，节奏用“长—短—短—长”。先只写这一小节，按 Space 听轮廓。
  \item 用“复制片段”（Ctrl+C）与“粘贴片段”（Ctrl+V）把这一小节复制到第 2 小节；选中第 2 小节的四个音，手动把每个音往上挪两行，得到 D–E–F–A 的模进。
  \item 第 3、4 小节写答句：前两拍接住动机，句尾落到 C4；问句的第 2 小节句尾落在 G4，先不要落到 C。
  \item 打开节拍器，用 Space 或循环播放（C）反复听：问句结束时耳朵是否还想听下一句？答句结束时是否真的落地？
  \item 做两次变奏：第一次只改节奏（把第一个音加倍），第二次只改轮廓（把上行改成倒影）；和弦音与句尾主音保持不变。
\end{enumerate}

\screenshot{piano-roll-four-bar-question-answer}{四小节问句与答句}

\textbf{范围说明：}夜半当前没有自动旋律生成、动机识别、乐句分析或“问句/答句”判定；这些只能靠耳朵与手动改音。和弦轨是固定五块演示数据，不跟随工程和弦；需要真实配和声时，请在卷帘里另写和弦轨或音频轨。应用也不会替你判断某个音是“张力音”还是“错音”，那取决于上下文与风格。

\subsection{练习}

\begin{enumerate}
  \item 以 C–D–E–G 为原型，写出上行二度的模进、以 C 为轴的倒影、以及逆行。
  \par\noindent 答案：模进 D–E–F–A；倒影 C–B–A–F；逆行 G–E–D–C。
  \item 写四小节问答句：问句两小节停在 G4，答句两小节回到 C4；说明你为什么这样收。
  \par\noindent 答案：不唯一。合格的标准是问句末音不是主音、答句末音是 C4，且答句能听出对问句动机的回应。
  \item C 大三和弦上，C、E、G、D、F、A、B 各是什么角色？
  \par\noindent 答案：C/E/G 和弦音；D 可作经过音或 9 音；F 可作辅助音、倚音或张力音；A 是 6 音张力；B 是大七张力/导音。
  \item 同样四个音，只改节奏与只改音高，哪一种更容易保住“动机身份”？为什么？
  \par\noindent 答案：通常节奏更容易被认出；音高轮廓也重要，两者一起变最容易丢掉原型。
  \item 在夜半把第 1 小节的四拍动机复制到第 2 小节，并做成模进；写出你用的菜单/快捷键。
  \par\noindent 答案：复制片段 Ctrl+C、粘贴片段 Ctrl+V；选中第 2 小节四个音后手动上移两行。
\end{enumerate}

\subsection{帮助主题映射}

\begin{center}
\small
\begin{tabular}{p{0.30\linewidth}p{0.42\linewidth}p{0.20\linewidth}}
\toprule
本章小节 & 应用内帮助锚点 & 备注 \\
\midrule
动机 & \texttt{music-theory-motive} & 概念 \\
乐句与问答 & \texttt{music-theory-phrase} & 四小节结构 \\
重复与模进 & \texttt{music-theory-sequence} & 表 7-1 \\
旋律与和弦 & \texttt{music-theory-melody-harmony} & 表 7-2 \\
节奏动机 & \texttt{music-theory-rhythm-motive} & 表 7-3 \\
在夜半里怎么做 & \texttt{music-theory-piano-roll-melody} & 操作 \\
练习 & \texttt{music-theory-ch7-practice} & 答案在章末 \\
\bottomrule
\end{tabular}
\end{center}
